% GENERATED FILE. Edit canonical lessons, then run npm run generate:books.
% canonical-source-hash: fnv1a64:22baa331797935fe
% canonical-lessons: GE-C512-schlitten, GE-C512-rodeln, GE-C512-schlittschuh, GE-C512-piste, GE-C512-skilift

\chapter{Im Schnee}
\label{ch:ge-im-schnee}
\hlchaptermodality{512}

\begin{chapteropening}
\textbf{By the end of this chapter:} \emph{I can talk about winter sports.}

Complete the last lesson of chapter 512: I can talk about winter sports.
\end{chapteropening}

\section[der Schlitten]{der Schlitten — sledge, sled}

\label{lesson:GE-C512-schlitten}



\begin{quote}
Before the new one: say the German for to paddle, then the German for canoe.
\end{quote}

\subsection*{You'll want to know: der Schlitten}
\textbf{der Schlitten} — \textquotedblleft{}sledge, sled\textquotedblright{}.

\begin{quote}
Die Kinder fahren mit dem Schlitten den Hügel hinunter.
\end{quote}

\begin{quote}
\emph{The children sledge down the hill.}
\end{quote}

Plural: \emph{die Schlitten}. \emph{Schlitten fahren} — to go sledging.

\begin{cousinweb}
\emph{Schlitten} is inherited from Middle High German \emph{slite} and Old High German \emph{slito} \textquotedblleft{}sled\textquotedblright{}, from Proto-Germanic \emph{slidô} (EtymDB); English \emph{sled} is a cousin.
\end{cousinweb}

\begin{grammarlens}[title={What you've built}]
One more word for telling, explaining and describing at B1.
\end{grammarlens}

\subsection*{Your turn}
\begin{itemize}
\raggedright
  \item \emph{Say it:} \emph{der Schlitten}
  \item \emph{Say it:} \emph{der Schlitten}, once more
  \item \emph{Say it:} a sentence of your own with \emph{der Schlitten}
  \item \emph{Recall:} say the German for to paddle, then the German for canoe, then say \emph{der Schlitten} again
\end{itemize}

\begin{tcolorbox}[breakable,colback=teal!4,colframe=teal!35!black,title={Before you move on}]
What does der Schlitten mean? (\textquotedblleft{}Sledge, sled\textquotedblright{}.) Say it once more.
\end{tcolorbox}
\section[rodeln]{rodeln — to toboggan}

\label{lesson:GE-C512-rodeln}



\begin{quote}
Before the new one: say the German for canoe, then the German for sledge.
\end{quote}

\subsection*{You'll want to know: rodeln}
\textbf{rodeln} — \textquotedblleft{}to toboggan\textquotedblright{}.

\begin{quote}
Am Sonntag sind wir im Wald gerodelt.
\end{quote}

\begin{quote}
\emph{On Sunday we went tobogganing in the woods.}
\end{quote}

Perfekt with \emph{sein} or \emph{haben}. Common in southern Germany and Austria; elsewhere also \emph{Schlitten fahren}.

\begin{cousinweb}
\emph{rodeln} is a word from the south of the German-speaking area; its origin is not settled, and EtymDB lists it without a source.
\end{cousinweb}

\begin{grammarlens}[title={What you've built}]
One more word for telling, explaining and describing at B1.
\end{grammarlens}

\subsection*{Your turn}
\begin{itemize}
\raggedright
  \item \emph{Say it:} \emph{rodeln}
  \item \emph{Say it:} \emph{rodeln}, once more
  \item \emph{Say it:} a sentence of your own with \emph{rodeln}
  \item \emph{Recall:} say the German for canoe, then the German for sledge, then say \emph{rodeln} again
\end{itemize}

\begin{tcolorbox}[breakable,colback=teal!4,colframe=teal!35!black,title={Before you move on}]
What does rodeln mean? (\textquotedblleft{}To toboggan\textquotedblright{}.) Say it once more.
\end{tcolorbox}
\section[der Schlittschuh]{der Schlittschuh — ice skate}

\label{lesson:GE-C512-schlittschuh}



\begin{quote}
Before the new one: say the German for sledge, then the German for to toboggan.
\end{quote}

\subsection*{You'll want to know: der Schlittschuh}
\textbf{der Schlittschuh} — \textquotedblleft{}ice skate\textquotedblright{}.

\begin{quote}
Im Winter laufen wir auf dem See Schlittschuh.
\end{quote}

\begin{quote}
\emph{In winter we go ice-skating on the lake.}
\end{quote}

Plural: \emph{die Schlittschuhe}. \emph{Schlittschuh laufen} — to skate.

\begin{cousinweb}
\emph{Schlittschuh} is built on a first part related to \emph{Schlitten} \textquotedblleft{}sledge\textquotedblright{} and \emph{Schuh} \textquotedblleft{}shoe\textquotedblright{}, inherited from Old High German \emph{scuoh}, from Proto-Germanic \emph{skōhaz} (EtymDB).
\end{cousinweb}

\begin{grammarlens}[title={What you've built}]
One more word for telling, explaining and describing at B1.
\end{grammarlens}

\subsection*{Your turn}
\begin{itemize}
\raggedright
  \item \emph{Say it:} \emph{der Schlittschuh}
  \item \emph{Say it:} \emph{der Schlittschuh}, once more
  \item \emph{Say it:} a sentence of your own with \emph{der Schlittschuh}
  \item \emph{Recall:} say the German for sledge, then the German for to toboggan, then say \emph{der Schlittschuh} again
\end{itemize}

\begin{tcolorbox}[breakable,colback=teal!4,colframe=teal!35!black,title={Before you move on}]
What does der Schlittschuh mean? (\textquotedblleft{}Ice skate\textquotedblright{}.) Say it once more.
\end{tcolorbox}
\section[die Piste]{die Piste — ski run, slope}

\label{lesson:GE-C512-piste}



\begin{quote}
Before the new one: say the German for to toboggan, then the German for ice skate.
\end{quote}

\subsection*{You'll want to know: die Piste}
\textbf{die Piste} — \textquotedblleft{}ski run, slope\textquotedblright{}.

\begin{quote}
Die Piste war voll mit Skifahrern.
\end{quote}

\begin{quote}
\emph{The slope was full of skiers.}
\end{quote}

Plural: \emph{die Pisten}. Also a runway.

\begin{cousinweb}
\emph{Piste} was borrowed from French \emph{piste} \textquotedblleft{}track, trail\textquotedblright{}, which came from Italian \emph{pista}; EtymDB has no entry for the German word.
\end{cousinweb}

\begin{grammarlens}[title={What you've built}]
One more word for telling, explaining and describing at B1.
\end{grammarlens}

\subsection*{Your turn}
\begin{itemize}
\raggedright
  \item \emph{Say it:} \emph{die Piste}
  \item \emph{Say it:} \emph{die Piste}, once more
  \item \emph{Say it:} a sentence of your own with \emph{die Piste}
  \item \emph{Recall:} say the German for to toboggan, then the German for ice skate, then say \emph{die Piste} again
\end{itemize}

\begin{tcolorbox}[breakable,colback=teal!4,colframe=teal!35!black,title={Before you move on}]
What does die Piste mean? (\textquotedblleft{}Ski run, slope\textquotedblright{}.) Say it once more.
\end{tcolorbox}
\section[der Skilift]{der Skilift — ski lift}

\label{lesson:GE-C512-skilift}



\begin{quote}
Before the new one: say the German for ice skate, then the German for ski run.
\end{quote}

\subsection*{You'll want to know: der Skilift}
\textbf{der Skilift} — \textquotedblleft{}ski lift\textquotedblright{}.

\begin{quote}
Am Skilift mussten wir lange warten.
\end{quote}

\begin{quote}
\emph{We had to wait a long time at the ski lift.}
\end{quote}

Plural: \emph{die Skilifte}. \emph{Ski} sounds like \emph{Schi}.

\begin{cousinweb}
\emph{Skilift} is built on \emph{Ski}, borrowed from Norwegian \emph{ski} \textquotedblleft{}a strip of wood\textquotedblright{}, and \emph{Lift}, borrowed from English \emph{lift}; EtymDB lists both German words without a source.
\end{cousinweb}

\begin{grammarlens}[title={What you've built}]
One more word for telling, explaining and describing at B1.
\end{grammarlens}

\subsection*{Your turn}
\begin{itemize}
\raggedright
  \item \emph{Say it:} \emph{der Skilift}
  \item \emph{Say it:} \emph{der Skilift}, once more
  \item \emph{Say it:} a sentence of your own with \emph{der Skilift}
  \item \emph{Recall:} say the German for ice skate, then the German for ski run, then say \emph{der Skilift} again
\end{itemize}

\begin{tcolorbox}[breakable,colback=teal!4,colframe=teal!35!black,title={Before you move on}]
What does der Skilift mean? (\textquotedblleft{}Ski lift\textquotedblright{}.) Say it once more.
\end{tcolorbox}
